
The one-sided communication semantics defined by the \openshmem specification
apply to symmetric objects allocated from the \ac{GPU} symmetric heap. All
routines defined in \openshmem 1.6---including \ac{RMA}, \ac{AMO}, signaling,
synchronization, and collective operations---retain their specified behavior and
may be used with GPU-attached memory buffers.

The remote operand in a communication operation may reside in either the default
(CPU-attached) symmetric heap defined by \openshmem 1.6 or the \ac{GPU}
symmetric heap for the GPU-aware communication operations. For GPU-centric
communication operations, the remote operand shall reside in the \ac{GPU}
symmetric heap only. The same \openshmem routine shall be used regardless of whether
the remote memory is CPU-attached or GPU-attached.

\subsection{GPU-Aware Communication}\label{subsec:gpu_aware_comm}

GPU-aware communication allows GPU-attached memory buffers to be used as
operands in existing \openshmem routines without modification to interfaces or
calling conventions. GPU-attached memory buffers may be used as:

\begin{itemize}
\item \textbf{Source operands} in put operations
\item \textbf{Destination operands} in get operations
\item \textbf{Local operands} in atomic memory operations
\item \textbf{Target buffers} in point-to-point synchronization operations
\item \textbf{Local buffers} in collective operations
\end{itemize}

When a GPU-attached memory buffer is used as an operand in an \openshmem
routine, the \openshmem implementation shall make the contents of that buffer
available to the communication operation and shall make updates visible according
to the completion and ordering semantics defined in \openshmem 1.6. This
requirement does not prescribe the implementation mechanism, such as direct
network access to GPU-attached memory, staging through CPU-attached memory, or
runtime synchronization with the \ac{GPU} memory hierarchy.

\subsection{GPU-Centric Communication}\label{subsec:gpu_centric_comm}

GPU-centric communication allows \ac{GPU} threads to invoke put, get, atomic,
signaling, synchronization, and collective operations directly from \ac{GPU}
execution contexts. Device-initiated
operations use the same addressing model defined by the \openshmem
specification: a symmetric address and a target \ac{PE} identifier. The
symmetric address shall be within the allocated region returned by the \ac{GPU} symmetric heap
allocation routine and is valid only at the \ac{PE} where the allocation was
performed.

\textbf{Note:} Support for GPU-centric communication is independent of support
for GPU-aware communication; an implementation may support either or both.

\subsection{Supported Operations}\label{subsec:supported_operations}

The following operations are supported on symmetric objects in the \ac{GPU}
symmetric heap:

\begin{longtable}{|p{5cm}|p{9cm}|}
\hline
\textbf{Operation} & \textbf{Description} \\
\hline
\endfirsthead
\hline
\textbf{Operation} & \textbf{Description} \\
\hline
\endhead
\ac{RMA} & Put and get operations, including bulk and single-element variants \\
\hline
\ac{AMO} & Fetch, non-fetch, and compare-and-swap variants \\
\hline
Signaling Operations & Combined data transfer and signal update at the target \ac{PE} \\
\hline
Synchronization & Barrier, fence, quiet, and point-to-point wait/test primitives \\
\hline
\end{longtable}

\subsection{Ordering and Completion}\label{subsec:ordering_completion}

Ordering and completion semantics follow the \openshmem memory model. Updates to
symmetric objects are unordered by default. A \ac{PE} may enforce point-to-point
ordering using \FUNC{shmem\_fence}. \FUNC{shmem\_quiet} may be used to ensure
completion of prior communication operations issued by the calling \ac{PE}.
Barrier and point-to-point wait/test operations may be used to coordinate access
across \acp{PE}.

